\begin{scanpage}{0065}{48}
Consequently, in the physical optics approach the field scattered in the directions $\theta=0$ and $\theta=\pi$ preserves the polarization of the incident wave.
\sectiontitle{8}{The Field From the Uniform Part of the Current}
Let us proceed to calculation of the field created by the nonuniform part of the current with normal irradiation of the disk. Since the latter is concentrated mainly in the vicinity of the disk's edge $(\rho=a)$, the vector potential corresponding to it will equal, in accordance with Equation (7.06),
\[\mathbf A=\frac ac\frac{e^{ikR}}R\int_0^a d\rho\int_0^{2\pi}\mathbf j^1(\rho,\psi)e^{-ik\rho\sin\theta\cos(\psi-\phi)}\,d\psi.\tag{8.01}\]
The inner integral is calculated with $ka\sin\theta\gg1$ based on the stationary phase method (see, for example [21], p.\ 256), and Equation (8.01) is transformed to the form
\[\begin{aligned}\mathbf A(x,y,z)&=\frac ac\sqrt{\frac{2\pi}{ka\sin\theta}}\frac{e^{ikR}}R e^{i\frac\pi4}\times\\&\quad\times\left[\int_0^a\mathbf j^1(\rho,\psi_1)e^{-ik\rho\sin\theta}\,d\rho-i\int_0^a\mathbf j^1(\rho,\psi_2)e^{ik\rho\sin\theta}\,d\rho\right],\end{aligned}\tag{8.02}\]
which allows one to interpret the fringing field as the field from a luminous line on the disk. This line is a diameter, the polar angle $\psi$ of the points on which equals
\[\psi_1=\phi\quad\text{and}\quad\psi_2=\pi+\phi.\tag{8.03}\]
Assuming the diameter of the disk is sufficiently large in comparison with the wavelength $(ka\gg1)$, one may approximately assume that the nonuniform part of the current near the disk's edge will be the same as on the corresponding half-plane (Figure 19). On the basis of \S\ 4, the field from the nonuniform part of the current flowing on the half-plane $-\infty<y_1<a$ may be represented in the form
\end{scanpage}
